% Generated draft: requires mathematical and source review.
\documentclass[11pt]{article}
\usepackage[margin=25mm]{geometry}
\usepackage{fontspec}
\setmainfont{Latin Modern Roman}
\usepackage{amsmath,amssymb,mathtools,unicode-math}
\setmathfont{Latin Modern Math}
\usepackage{enumitem,needspace,xurl,hyperref}
\setlength{\parindent}{0pt}
\setlength{\parskip}{5pt}
\setlength{\emergencystretch}{3em}
\newcommand{\N}{\mathbb{N}}
\newcommand{\Z}{\mathbb{Z}}
\newcommand{\Q}{\mathbb{Q}}
\newcommand{\R}{\mathbb{R}}
\newcommand{\C}{\mathbb{C}}
\newcommand{\F}{\mathbb{F}}
\newcommand{\A}{\mathbb{A}}
\newcommand{\PP}{\mathbb{P}}
\newcommand{\E}{\mathbb{E}}
\newcommand{\eps}{\varepsilon}
\newcommand{\dd}{\,\mathrm{d}}
\begin{document}
\section*{UnsolvedMath ID: 30003776}
\textbf{Source problem code:} catalogue entry
\subsection{Definitions and mathematical statement}
One standard framework for the question is the following; the supplied record does not explicitly fix these conventions. An expander sequence is a sequence \((G_n)_{n\in\N}\) of finite connected undirected simple graphs, with \(|V(G_n)|\geq 2\) and \(|V(G_n)|\to\infty\), for which there exist constants \(D<\infty\) and \(h>0\) such that, for every \(n\),
\[
\max_{v\in V(G_n)}\deg(v)\leq D,
\qquad
\min_{\substack{\varnothing\ne A\subseteq V(G_n)\\ |A|\leq |V(G_n)|/2}}
\frac{|\partial_E A|}{|A|}\geq h.
\]
Here \(\partial_E A\) is the set of edges with exactly one endpoint in \(A\), and each graph carries its shortest-path metric. These are conventional working definitions, not hypotheses explicitly documented in the supplied problem statement.

For metric spaces \((X,d_X)\) and \((Y,d_Y)\), a map \(f:X\to Y\) is bornologous if
\[
\forall R\geq 0\ \exists S\geq 0\ \forall x,x'\in X,
\qquad d_X(x,x')\leq R\ \Longrightarrow\ d_Y(f(x),f(x'))\leq S.
\]
A standard meaning of coarse equivalence is the existence of bornologous maps \(f:X\to Y\) and \(g:Y\to X\) and a constant \(C<\infty\) satisfying
\[
 d_X(g(f(x)),x)\leq C\quad(x\in X),
 \qquad
 d_Y(f(g(y)),y)\leq C\quad(y\in Y).
\]
This supplies a reference definition, but the source asks which notion is appropriate and does not prescribe this relation as the desired classification relation.

\textbf{Research question.} Identify and justify a notion of coarse equivalence or coarse distinction that gives a useful classification of expander sequences. The supplied record does not specify whether the objects compared are coarse disjoint unions, sequences compared through uniformly controlled maps between components, or another construction. It also does not specify the meaning of a useful classification, the permitted invariants, or the required completeness of such invariants. Accordingly, the source supports an open-ended classification problem, not a uniquely determined conjecture with fixed quantifiers and a specified conclusion.

The original wording also asked whether all expanders are superexpanders. The supplied correction explicitly removes that clause because it duplicates source problem \texttt{OWR-16160-001}; it is not an additional target of this entry.
\subsection{Short English statement}
Find an appropriate way to compare or distinguish expander sequences by their large-scale geometry and use it to organize them into meaningful classes. The source leaves the comparison relation and the requirements for a classification unspecified. Its corrected formulation excludes the separate question of whether every expander is a superexpander.
\subsection{Sources}
\begin{itemize}
\item Local UnsolvedMath problem registry, permanent ID \(30003776\), source code \texttt{catalogue entry}, in the Oberwolfach Reports Open Problems collection. The supplied record gives proposed year \(2018\), names no individual proposer, and labels the problem \texttt{open}. Its statement status is \texttt{corrected\_verified}; the accompanying registry explanation says that the all-expanders clause duplicated \texttt{OWR-16160-001} and was removed. These are supplied curation labels, not independent verification.
\item Supplied report citation: \emph{Mini-Workshop: Superexpanders and Their Coarse Geometry (2019)}, Oberwolfach Reports, pp.\ 1117--1160, DOI \texttt{10.4171/owr/2018/19}. \url{https://doi.org/10.4171/owr/2018/19}. The background attributes the problem to workshop open-problem discussions documented in this report; no precise internal problem location is supplied.
\item Supplied literature source: Mendel--Naor, nonlinear spectral calculus. \url{https://arxiv.org/abs/1207.4705}.
\item Supplied literature source: OWR 2018 report. \url{https://doi.org/10.4171/owr/2018/19}.
\item Literature assessment checked at \texttt{2026-08-21}, according to the supplied record. This is supplied dated curation derived from the OWR material, not independently verified here. It labels the problem open and reports that known constructions and coarse invariants provide many superexpanders and distinguish some expander families, while neither the all-expanders conjecture nor a general correct coarse-equivalence classification is settled. The supplied background explicitly requests independent verification of time-sensitive status and cited claims.
\end{itemize}
\end{document}
